\documentclass[11pt]{article}
\usepackage{amsmath,amssymb,amsthm}
\newtheorem{proposition}{Proposition}
\newtheorem{lemma}[proposition]{Lemma}
\newcommand{\Z}{\mathbb Z}
\newcommand{\Q}{\mathbb Q}
\begin{document}
\title{Formal ratios with one free elliptic index}
\author{}\date{10 November 2026}\maketitle

These statements concern specified replacements of the current elliptic
source. They give no lower bound for all Diophantine definitions of the
integers. In particular, they do not exclude a different arithmetic
source, a different prime construction, or correspondences not covered
by their hypotheses.

\begin{lemma}
Let $E$ be an elliptic curve over a number field and let $P$ be a
nontorsion rational point. For each positive integer $B$, there are
only finitely many good finite places at which the reduction order of
$P$ is at most $B$.
\end{lemma}
\begin{proof}
At such a place, $[m]P$ reduces to the identity for some
$1\le m\le B$. Every $[m]P$ is a fixed nonidentity point. Its
reduction is the identity at only finitely many places, as is seen
from an integral Weierstrass equation and the denominator of its
abscissa. The union over the finitely many $m$ is finite.
\end{proof}

\begin{proposition}
Fix integer coefficients $a_i,b_i,c_i$ for $i=0,1$, and set
\[
        L_i(n,h)=a_i n+b_i h+c_i.
\]
Let $I\subset\Z$ be infinite. Suppose that for every $n\in I$
there are unbounded positive integers $M$ and integers $h$ for which
\[
                  M\mid L_0(n,h),\qquad M\mid L_1(n,h).
\]
Then
\[
        b_1a_0-b_0a_1=0,
        \qquad b_1c_0-b_0c_1=0.
\]
If $b_0\ne0$, the two affine forms are proportional over $\Q$,
with $L_1=(b_1/b_0)L_0$. If $b_0=0$, then $L_0$ is identically
zero.
\end{proposition}
\begin{proof}
Eliminating $h$ gives
\[
 M\mid b_1L_0-b_0L_1
      =(b_1a_0-b_0a_1)n+b_1c_0-b_0c_1.
\]
For fixed $n$, unboundedness of $M$ forces the integer on the
right to be zero. It vanishes at infinitely many $n$, so its two
coefficients vanish. This proves proportionality if $b_0\ne0$.
If $b_0=0$, the original divisibility $M\mid a_0n+c_0$ and
unboundedness of $M$ force $a_0n+c_0=0$ for every $n\in I$;
hence $a_0=c_0=0$.
\end{proof}

Applied to a fixed root point $[n]P$ and one free auxiliary point
$[h]P$, the proposition excludes a varying formal-logarithm ratio
whose numerator and denominator are fixed affine combinations of
those points, if both must be formal at arbitrarily large comparison
primes. The denominator form cannot be identically zero. Otherwise
the exact ratio is the fixed rational number $b_1/b_0$. This explains
why simply replacing the third point by an addition formula does not
preserve the current completeness argument.

\begin{proposition}
Let $F\in\Q[H,U]$, let $d>0$ satisfy $dF\in\Z[H,U]$, and
fix a rational number $\beta$. Suppose that unbounded positive
integers $M$ admit integer pairs $(h,u)$ satisfying
\[
           M\mid h,\quad M\mid u,\quad F(h,u)=\beta.
\]
Then $\beta=F(0,0)$.
\end{proposition}
\begin{proof}
The last equation makes $d\beta$ an integer and gives
$M\mid d\beta-dF(0,0)$. Unboundedness of $M$ proves the claim.
\end{proof}

This excludes a second immediate replacement: two free primary-formal
points whose index polynomial is asserted to equal the input. It does
not exclude rational functions of the indices on a properly constrained
subset. For example, the rational function
\[
       \frac{u-1}{h-1}-\frac uh
          =\frac{u-h}{h(h-1)}
\]
can equal every rational input by choosing $h$ divisible by its
denominator. Its existence alone therefore supplies no integrality
criterion.

\begin{proposition}
Let $E$ be an elliptic curve in characteristic zero, let $f$ be a
nonconstant rational function on $E$, and let $m\in\Z$ have
$|m|>1$. There are only finitely many points $Q$ in the algebraic
closure at which $f([m]Q)=mf(Q)$ is defined and holds.
\end{proposition}
\begin{proof}
If there were infinitely many, the two rational functions would be
identical. But the degree of $f\circ[m]$ is $m^2\deg f$, whereas
the degree of $mf$ is $\deg f$, a contradiction.
\end{proof}

Thus a rational formal parameter such as $-x/y$ cannot replace the
elliptic logarithm by an exactly linear function on infinitely many
auxiliary multiples for a fixed input. In particular, imposing the
exact monic Lucas equation directly on its ratios has a rational
completeness requirement which is not supplied by the current source.

\section*{The monic Lucas possibility}

For integer $a,b$, a rational solution $t$ of
\[
 t^2-a(a-2)tb+(a-2)^2b^2-1=0
\]
is an integer, by the rational-root theorem. At the level of integer
labels one can therefore replace the $t$ label by a finite rational
root. The existing source, however, exposes these labels only through
truncated local ratios. Applying this monic equation to those ratios
does not give its rational completeness. Adding separate rational
coordinates for $a,b$ instead requires a new height argument proving
that these coordinates equal the elliptic indices. The existing
single trace scale does not control their arbitrary denominators.
No coordinate reduction is inferred from monicity alone.
\end{document}
